\documentclass[10pt,a4paper]{amsart}
\usepackage[margin=19mm]{geometry}
\usepackage{amsmath,amssymb,mathtools}
\usepackage[scr=rsfs,cal=euler]{mathalfa}
\usepackage{xfrac}
\usepackage[unicode,hidelinks,bookmarks=false]{hyperref}
\newcommand{\ie}{i.e.\ }
\newcommand{\cf}{cf.\ }
\newcommand{\resp}{respectively\ }
\newcommand{\Subsection}{\subsection}
\providecommand{\nobreakdash}{\nobreak\hbox{-}}
\setlength{\emergencystretch}{2em}
\multlinegap0pt
\usepackage{amssymb}
\newtheorem{lemma}{Lemma}
\newtheorem{proposition}{Proposition}
\title{An Exhaustive Census of Main-Diagonal Symmetric Costas Arrays of Orders 37--42}
\author{Bogdan Dumitru, Cristian Budala}
\date{Source: arXiv:2608.28690v1 (CC BY 4.0)}
\begin{document}
\maketitle

\noindent\emph{Source and licence.} An Exhaustive Census of Main-Diagonal Symmetric Costas Arrays of Orders 37--42. Version arXiv:2608.28690v1 (CC BY 4.0), distributed by arXiv under the Creative Commons Attribution 4.0 International licence, http://creativecommons.org/licenses/by/4.0/, as recorded in the arXiv licence metadata for this exact version and shown on the abstract page https://arxiv.org/abs/2608.28690v1. Full licence notice: \texttt{licenses/arxiv-2608.28690-CC-BY-4.0.txt}.\par\smallskip

\noindent\textit{Editorial scope.} Counted unit: the article's proposition prop:bulk (source0.tex lines 379--385). The article's complete original proof of it is delivered with it (source0.tex lines 387--498, one \texttt{proof} environment of 112 lines). No mathematics is written, repaired or re-derived: the statement and the proof are the article's own lines, in the article's own notation, and no retained line is substituted or re-flowed.  Retained uncounted as the source-local context the proof needs: lines 322--330.  Nothing was omitted from any retained span, so the package carries no omission record.  No retained line contains a citation, so the wrapper carries no bibliography and no bibliography entry.  No retained line contains a figure environment, an image-inclusion command or drawing code, so no image is delivered and the package declares no asset dependency.  Editorial formatting only, in addition: an AMS \texttt{amsart} preamble that restates the article's own declarations and macros used by the retained lines, namely the environments \texttt{lemma}, \texttt{proposition} on separate counters, together with the standard packages those lines need.  The retained environments print in this document as Lemma 1, Proposition 1 in order of appearance, whereas the article prints the same items with the numbering of its own full document.  The complete native source of the article as distributed in the arXiv e-print for this exact version is bundled unchanged as \texttt{sources/208802/source0.tex}.
\begin{lemma}
\label{lem:involution-ratio}
Let $I_n$ be the number of involutions on $n$ labeled points.  For every
fixed nonnegative integer $r$,
\[
\frac{I_{n-r}}{I_n}
=n^{-r/2}\left(1-\frac{r}{2\sqrt n}+O(n^{-1})\right).
\]
\end{lemma}

\begin{proposition}
\label{prop:bulk}
\[
\mathbb{E}[W_{\mathrm{full}}]
=\frac{n^2}{9}+\frac{n^{3/2}}{180}+O(n).
\]
\end{proposition}

\begin{proof}
For a fixed stride $k$, a collision comparison has the form
\[
p(i+k)-p(i)=p(j+k)-p(j),\qquad i<j.
\tag{1}
\]
There are
$\sum_{k=1}^{n-1}\binom{n-k}{2}=\binom n3$ such comparisons.  The cases
$j=i+k$, in which only three row positions occur, number $O(n^2)$ and
will be included in the remainder.

\emph{Eight vertices.}
First suppose that the four row positions and their four images are
disjoint.  For each row comparison, the unrestricted number of image
quadruples satisfying (1) is
\[
\sum_{d\ne0}(n-|d|)^2
=\frac{n(n-1)(2n-1)}{3}
=\frac23n^3+O(n^2).
\]
Requiring the images to be distinct and disjoint from the rows removes
only $O(n^2)$ quadruples.  Hence the number $A_8(n)$ of compatible
eight-vertex specifications is
\[
A_8(n)=\frac19n^6+O(n^5).
\]
Every such partial involution extends in $I_{n-8}$ ways, so
Lemma~\ref{lem:involution-ratio} gives
\[
A_8(n)\frac{I_{n-8}}{I_n}
=\frac19n^2-\frac49n^{3/2}+O(n).
\tag{2}
\]

\emph{Seven vertices.}
A compatible specification on exactly seven vertices has one fixed row
position and pairs the other three rows with external vertices.  If the
fixed position is $x$, first ignore assignments that reuse a row or an
image.  The number of assignments of the other three images is
\[
F_n(x)=\sum_{d=-x}^{n-1-x}(n-|d|)
=\frac{n^2}{2}+nx+\frac n2-x^2-x.
\tag{3}
\]
We now sum this polynomial over
$1\leq k\leq n-1$ and $0\leq i<j\leq n-k-1$.  The standard sums of
powers give
\begin{align*}
\sum_{k,i,j}F_n(i)
=\sum_{k,i,j}F_n(j+k)
&=\frac{n(n-1)(n-2)(13n^2-n+6)}{120},\\
\sum_{k,i,j}F_n(i+k)
=\sum_{k,i,j}F_n(j)
&=\frac{n(n-1)(n-2)(7n^2+n+4)}{60}.
\end{align*}
Thus the four possible fixed positions give
\[
\frac{n(n-1)(n-2)(27n^2+n+14)}{60}
=\frac9{20}n^5+O(n^4).
\]
For each row comparison, forbidding a reused row or image removes only
$O(n)$ assignments; the comparisons with an overlapping row contribute
only $O(n^4)$ to the displayed unrestricted sum.  Therefore the number
of compatible seven-vertex specifications is
\[
A_7(n)=\frac9{20}n^5+O(n^4).
\]
Lemma~\ref{lem:involution-ratio} now gives
\[
A_7(n)\frac{I_{n-7}}{I_n}
=\frac9{20}n^{3/2}+O(n).
\tag{4}
\]

\emph{Remaining patterns.}
For four distinct row positions, let $f$ be the number that are fixed,
$t$ the number of internal transpositions among them, and $e$ the number
paired with external vertices.  Then
\[
f+2t+e=4,
\qquad\text{and the partial involution uses }4+e\text{ vertices}.
\]
The cases $e=4$ and $e=3$ are (2) and (4).  If $0<e\leq2$, equation (1)
is a nonzero affine equation in the $e$ external labels.  There are
$O(n^{e-1})$ label assignments for each row comparison, and the total
contribution is
\[
O\!\left(n^3 n^{e-1} n^{-(4+e)/2}\right)=O(n^{e/2})=O(n).
\]
For $e=0$, there are $O(n^3)$ specifications, each with probability
$O(n^{-2})$, again contributing $O(n)$.

When $j=i+k$, only three row positions occur and there are $O(n^2)$ row
comparisons.  If $e>0$ of the three rows are paired externally, the
arithmetic-progression equation among their images leaves
$O(n^{e-1})$ external assignments.  Since the partial involution uses
$3+e$ vertices and $e\leq3$, the contribution is
\[
O\!\left(n^2n^{e-1}n^{-(3+e)/2}\right)=O(n^{(e-1)/2})=O(n).
\]
The case $e=0$ contributes
$O(n^2n^{-3/2})=O(n^{1/2})$.

Combining (2) and (4), with all other patterns absorbed into $O(n)$,
gives
\[
\mathbb{E}[W_{\mathrm{full}}]
=\frac19n^2+
\left(-\frac49+\frac9{20}\right)n^{3/2}+O(n)
=\frac19n^2+\frac1{180}n^{3/2}+O(n).
\]
\end{proof}

\end{document}
